\section{问题重述}

\subsection{问题背景}

山区洪涝灾害具有突发性强、影响范围分散、道路供电通信设施同时受损等特点。连续强降雨引发的山洪、滑坡与道路塌方，会使部分居民点在短时间内与外界失去稳定的地面交通联系。2026 年 7 月，受台风“美莎克”带来的持续强降雨影响，广西横州市镇龙乡遭遇严重洪涝灾害，一度出现断路、断电、断网，全乡约 8000 人受困\cite{ref_xinhua1,ref_xinhua2}，食品、饮用水和药品需依靠徒步或空中投送进入受灾区域。无人机受地面道路条件影响较小，可作为应急运输的重要补充手段；但完成物资投送不仅取决于能否飞抵目标，还受载荷、水平飞行距离、爬升与下降高度、沿线地形净空、能源储备、任务时限与设备周转等多种因素制约。同时，运输无人机在爬升、巡航、下降及物资交接全过程需持续保持指挥、遥测与交接确认链路；当固定通信设施受损或山体遮挡导致地面网关无法覆盖时，还需空中中继无人机在适当位置和时段提供临时通信保障。

题目以镇龙乡及周边山区为地理背景设置\textbf{标准测试场景}：提供 1 个临时调度中心 $O01$、15 个服务区 $S_1\!\sim\!S_{15}$ 和 80 个不可拆货箱，设置 A/B/C 三种运输机型共 8 架实体运输无人机，并配套共享电池组、30\,m DEM 数据、中继无人机、可更换能源组件及无线链路参数等信息（设备参数量级参考工业级无人机与通信模块公开规格\cite{ref_dji,ref_doodle}）。本题研究无人机物资运输、通信中继与救援资源配置之间的协同关系。

\subsection{需要解决的问题}

\begin{itemize}[leftmargin=2em]
  \item \textbf{问题一（单点往返运输能力与货箱组批）}：不考虑实体无人机和共享电池调度，每个架次采用 $O01\to S_i\to O01$ 的直接往返、仅服务一个服务区、不得跨服务区组批。需（1）计算三种机型在不同服务区执行单点往返任务时的最大安全载荷；（2）在货箱不可拆分、每箱只安排一次并满足载质量、装载体积与返航安全能量余量的条件下确定各服务区货箱组批方案；（3）综合往返架次数、总运输能耗与累计作业时间优化组批方案并说明各指标权衡关系；（4）讨论返航安全余量变化对最大安全载荷与组批结果的影响。
  \item \textbf{问题二（异构无人机多点多架次运输调度）}：不考虑通信保障，每个架次可访问一个或多个服务区并返回 $O01$。在满足货箱不可拆分、载质量与体积、返航安全余量、无人机与共享电池数量、充电周转及物资时限（医疗物资期望送达、首批保障截止）等约束下，联合确定货箱组批、访问顺序、机型、执行无人机、共享电池及各架次开始时刻，综合配送及时性、完工时间、运输能耗与架次数进行优化，并给出资源使用与可行性检验。
  \item \textbf{问题三（通信约束下的运输与中继联合调度）}：运输无人机在爬升、巡航、下降及投送阶段均应保持连续通信，直连不可用时可由一架中继无人机提供“运输—中继—$G01$”两段链路，通信状态按附录 3 判定。在满足时限、载荷与能量、连续通信及资源可用性约束下，联合确定运输方案与中继无人机的悬停位置、飞行海拔、服务时段及能源组件分配，综合及时性、联合完工时间、运输与中继总能耗及两类无人机架次进行优化。
  \item \textbf{问题四（救援任务分区与资源配置优化）}：以问题三方案为基础，将 15 个服务区分别划分为 2 个和 3 个任务组，保持问题三的组批、访问顺序、运输与中继任务及通信保障关系不变（多站架次须整组同划），各组独立执行、资源不得跨组调配，核算所需运输无人机、共享电池、中继无人机与中继能源组件数量，从资源规模、冗余、组间工作量均衡及与现有库存的缺口等方面比较两种分区方式，并对超库存情形给出缺口及原因。
\end{itemize}

各问遵循附录 2（时间、能源与充电周转规则）与附录 3（通信链路计算与服务规则）统一口径，公共物理规则、单位、对象编号和计算口径在四问间保持一致，并逐层继承数据。
